%% ODER: format ==         = "\mathrel{==}"
%% ODER: format /=         = "\neq "
%
%
\makeatletter
\@ifundefined{lhs2tex.lhs2tex.sty.read}%
  {\@namedef{lhs2tex.lhs2tex.sty.read}{}%
   \newcommand\SkipToFmtEnd{}%
   \newcommand\EndFmtInput{}%
   \long\def\SkipToFmtEnd#1\EndFmtInput{}%
  }\SkipToFmtEnd

\newcommand\ReadOnlyOnce[1]{\@ifundefined{#1}{\@namedef{#1}{}}\SkipToFmtEnd}
\usepackage{amstext}
\usepackage{amssymb}
\usepackage{stmaryrd}
\DeclareFontFamily{OT1}{cmtex}{}
\DeclareFontShape{OT1}{cmtex}{m}{n}
  {<5><6><7><8>cmtex8
   <9>cmtex9
   <10><10.95><12><14.4><17.28><20.74><24.88>cmtex10}{}
\DeclareFontShape{OT1}{cmtex}{m}{it}
  {<-> ssub * cmtt/m/it}{}
\newcommand{\texfamily}{\fontfamily{cmtex}\selectfont}
\DeclareFontShape{OT1}{cmtt}{bx}{n}
  {<5><6><7><8>cmtt8
   <9>cmbtt9
   <10><10.95><12><14.4><17.28><20.74><24.88>cmbtt10}{}
\DeclareFontShape{OT1}{cmtex}{bx}{n}
  {<-> ssub * cmtt/bx/n}{}
\newcommand{\tex}[1]{\text{\texfamily#1}}	% NEU

\newcommand{\Sp}{\hskip.33334em\relax}


\newcommand{\Conid}[1]{\mathit{#1}}
\newcommand{\Varid}[1]{\mathit{#1}}
\newcommand{\anonymous}{\kern0.06em \vbox{\hrule\@width.5em}}
\newcommand{\plus}{\mathbin{+\!\!\!+}}
\newcommand{\bind}{\mathbin{>\!\!\!>\mkern-6.7mu=}}
\newcommand{\rbind}{\mathbin{=\mkern-6.7mu<\!\!\!<}}% suggested by Neil Mitchell
\newcommand{\sequ}{\mathbin{>\!\!\!>}}
\renewcommand{\leq}{\leqslant}
\renewcommand{\geq}{\geqslant}
\usepackage{polytable}

%mathindent has to be defined
\@ifundefined{mathindent}%
  {\newdimen\mathindent\mathindent\leftmargini}%
  {}%

\def\resethooks{%
  \global\let\SaveRestoreHook\empty
  \global\let\ColumnHook\empty}
\newcommand*{\savecolumns}[1][default]%
  {\g@addto@macro\SaveRestoreHook{\savecolumns[#1]}}
\newcommand*{\restorecolumns}[1][default]%
  {\g@addto@macro\SaveRestoreHook{\restorecolumns[#1]}}
\newcommand*{\aligncolumn}[2]%
  {\g@addto@macro\ColumnHook{\column{#1}{#2}}}

\resethooks

\newcommand{\onelinecommentchars}{\quad-{}- }
\newcommand{\commentbeginchars}{\enskip\{-}
\newcommand{\commentendchars}{-\}\enskip}

\newcommand{\visiblecomments}{%
  \let\onelinecomment=\onelinecommentchars
  \let\commentbegin=\commentbeginchars
  \let\commentend=\commentendchars}

\newcommand{\invisiblecomments}{%
  \let\onelinecomment=\empty
  \let\commentbegin=\empty
  \let\commentend=\empty}

\visiblecomments

\newlength{\blanklineskip}
\setlength{\blanklineskip}{0.66084ex}

\newcommand{\hsindent}[1]{\quad}% default is fixed indentation
\let\hspre\empty
\let\hspost\empty
\newcommand{\NB}{\textbf{NB}}
\newcommand{\Todo}[1]{$\langle$\textbf{To do:}~#1$\rangle$}

\EndFmtInput
\makeatother
%
%
%
%
%
%
% This package provides two environments suitable to take the place
% of hscode, called "plainhscode" and "arrayhscode". 
%
% The plain environment surrounds each code block by vertical space,
% and it uses \abovedisplayskip and \belowdisplayskip to get spacing
% similar to formulas. Note that if these dimensions are changed,
% the spacing around displayed math formulas changes as well.
% All code is indented using \leftskip.
%
% Changed 19.08.2004 to reflect changes in colorcode. Should work with
% CodeGroup.sty.
%
\ReadOnlyOnce{polycode.fmt}%
\makeatletter

\newcommand{\hsnewpar}[1]%
  {{\parskip=0pt\parindent=0pt\par\vskip #1\noindent}}

% can be used, for instance, to redefine the code size, by setting the
% command to \small or something alike
\newcommand{\hscodestyle}{}

% The command \sethscode can be used to switch the code formatting
% behaviour by mapping the hscode environment in the subst directive
% to a new LaTeX environment.

\newcommand{\sethscode}[1]%
  {\expandafter\let\expandafter\hscode\csname #1\endcsname
   \expandafter\let\expandafter\endhscode\csname end#1\endcsname}

% "compatibility" mode restores the non-polycode.fmt layout.

\newenvironment{compathscode}%
  {\par\noindent
   \advance\leftskip\mathindent
   \hscodestyle
   \let\\=\@normalcr
   \let\hspre\(\let\hspost\)%
   \pboxed}%
  {\endpboxed\)%
   \par\noindent
   \ignorespacesafterend}

\newcommand{\compaths}{\sethscode{compathscode}}

% "plain" mode is the proposed default.
% It should now work with \centering.
% This required some changes. The old version
% is still available for reference as oldplainhscode.

\newenvironment{plainhscode}%
  {\hsnewpar\abovedisplayskip
   \advance\leftskip\mathindent
   \hscodestyle
   \let\hspre\(\let\hspost\)%
   \pboxed}%
  {\endpboxed%
   \hsnewpar\belowdisplayskip
   \ignorespacesafterend}

\newenvironment{oldplainhscode}%
  {\hsnewpar\abovedisplayskip
   \advance\leftskip\mathindent
   \hscodestyle
   \let\\=\@normalcr
   \(\pboxed}%
  {\endpboxed\)%
   \hsnewpar\belowdisplayskip
   \ignorespacesafterend}

% Here, we make plainhscode the default environment.

\newcommand{\plainhs}{\sethscode{plainhscode}}
\newcommand{\oldplainhs}{\sethscode{oldplainhscode}}
\plainhs

% The arrayhscode is like plain, but makes use of polytable's
% parray environment which disallows page breaks in code blocks.

\newenvironment{arrayhscode}%
  {\hsnewpar\abovedisplayskip
   \advance\leftskip\mathindent
   \hscodestyle
   \let\\=\@normalcr
   \(\parray}%
  {\endparray\)%
   \hsnewpar\belowdisplayskip
   \ignorespacesafterend}

\newcommand{\arrayhs}{\sethscode{arrayhscode}}

% The mathhscode environment also makes use of polytable's parray 
% environment. It is supposed to be used only inside math mode 
% (I used it to typeset the type rules in my thesis).

\newenvironment{mathhscode}%
  {\parray}{\endparray}

\newcommand{\mathhs}{\sethscode{mathhscode}}

% texths is similar to mathhs, but works in text mode.

\newenvironment{texthscode}%
  {\(\parray}{\endparray\)}

\newcommand{\texths}{\sethscode{texthscode}}

% The framed environment places code in a framed box.

\def\codeframewidth{\arrayrulewidth}
\RequirePackage{calc}

\newenvironment{framedhscode}%
  {\parskip=\abovedisplayskip\par\noindent
   \hscodestyle
   \arrayrulewidth=\codeframewidth
   \tabular{@{}|p{\linewidth-2\arraycolsep-2\arrayrulewidth-2pt}|@{}}%
   \hline\framedhslinecorrect\\{-1.5ex}%
   \let\endoflinesave=\\
   \let\\=\@normalcr
   \(\pboxed}%
  {\endpboxed\)%
   \framedhslinecorrect\endoflinesave{.5ex}\hline
   \endtabular
   \parskip=\belowdisplayskip\par\noindent
   \ignorespacesafterend}

\newcommand{\framedhslinecorrect}[2]%
  {#1[#2]}

\newcommand{\framedhs}{\sethscode{framedhscode}}

% The inlinehscode environment is an experimental environment
% that can be used to typeset displayed code inline.

\newenvironment{inlinehscode}%
  {\(\def\column##1##2{}%
   \let\>\undefined\let\<\undefined\let\\\undefined
   \newcommand\>[1][]{}\newcommand\<[1][]{}\newcommand\\[1][]{}%
   \def\fromto##1##2##3{##3}%
   \def\nextline{}}{\) }%

\newcommand{\inlinehs}{\sethscode{inlinehscode}}

% The joincode environment is a separate environment that
% can be used to surround and thereby connect multiple code
% blocks.

\newenvironment{joincode}%
  {\let\orighscode=\hscode
   \let\origendhscode=\endhscode
   \def\endhscode{\def\hscode{\endgroup\def\@currenvir{hscode}\\}\begingroup}
   %\let\SaveRestoreHook=\empty
   %\let\ColumnHook=\empty
   %\let\resethooks=\empty
   \orighscode\def\hscode{\endgroup\def\@currenvir{hscode}}}%
  {\origendhscode
   \global\let\hscode=\orighscode
   \global\let\endhscode=\origendhscode}%

\makeatother
\EndFmtInput
%
%
%
% First, let's redefine the forall, and the dot.
%
%
% This is made in such a way that after a forall, the next
% dot will be printed as a period, otherwise the formatting
% of `comp_` is used. By redefining `comp_`, as suitable
% composition operator can be chosen. Similarly, period_
% is used for the period.
%
\ReadOnlyOnce{forall.fmt}%
\makeatletter

% The HaskellResetHook is a list to which things can
% be added that reset the Haskell state to the beginning.
% This is to recover from states where the hacked intelligence
% is not sufficient.

\let\HaskellResetHook\empty
\newcommand*{\AtHaskellReset}[1]{%
  \g@addto@macro\HaskellResetHook{#1}}
\newcommand*{\HaskellReset}{\HaskellResetHook}

\global\let\hsforallread\empty

\newcommand\hsforall{\global\let\hsdot=\hsperiodonce}
\newcommand*\hsperiodonce[2]{#2\global\let\hsdot=\hscompose}
\newcommand*\hscompose[2]{#1}

\AtHaskellReset{\global\let\hsdot=\hscompose}

% In the beginning, we should reset Haskell once.
\HaskellReset

\makeatother
\EndFmtInput
% ----------------------
%% Stuff for TimCore


%% 'g' inferences





%% Intersection of effects

%% --------- Effects ------------


%% ----------------------


%% --------------------------
%%        Desugaring rules
%% --------------------------



% ----------------------------------------------------
% New desugaring, just does the wrapping part:
% ----------------------------------------------------
%   Input for wrapping


%%format Hide (e) = "{\mathcal H}\llbracket{}" e "\rrbracket{}"

%% 'si' is short for "solve or infer" or "unification variable or not" resp

%% 'md' is short for "mode" or "unification variable or not" resp

%% flipsi is a no-op for now


%% EXX is just EX, but with no argument




%% ----------- Applications --------------
%% %format applyx a b = a "@" b
%% %format applyx a b = a "\," b

%% Function application with no arguments

%% Function application with 1 argument

%% Function application with 1 argument, no parens

%% Function application with 1 argument, without parenthesis

%% Function application with many arguments

%% ----------- End of Applications --------------





%% Sequences e1; ...; en; v

%% xs etc stole syntax used in examples.

%% Old version: | for BNF, [] for choice
%%%format choice = "[ \! ]"
%%%format `choice` = "\mathop{[ \! ]}"
%% 1mm is about 0.3em with the current font

%% New version: tall | for BNF, fat | for choice
%%format vbar = "\mathop{\bigm|}"
%%format choice = "\boldsymbol{|}"
%%format `choice` = "\mathop{\boldsymbol{|}}"




%% %format ok = "\textbf{ok}"

%% %format `eqsemi` = "\hspace{-0.5mm}\fatsemi{}"
%% %format eqsemiKW = "\fatsemi"

%% Arrays and tuples

%% %format vmap (x) = "\llangle" x "\rrangle"





%%  OLD SYNTAX   %format split (x) (y) (z) = "\textbf{split}(" x ")\{" y "," z "\}"
%% %format defKW = "\textbf{def}"
% only used in comments, but still prettier this way:
%%%%format map = "\textbf{map}"

%% -------Functions ---------------


%% functionOC has open/closed subscript


%% -------End of functions ---------------





%% Unification
%%%format == = "\!=\!"


% format := = "\mapsto"
% %format squiggt = "\leadsto"
% %format squiggt = "\!-\!\!>\!\!\!"





























%% ----------- DPS Verse -------------

%% Translation from Mini Verse to DPS Verse


%% ----------- End of DPS Verse -------------






%% We combined S and C into one


%% ----------- Metatheory --------------


%% ----------------------------
%%         Verification
%% ----------------------------

%% assert = succeed, abbreviated to suc



%% Effects checking, like succeeds{e}





%% <fx> brackets

% Typesetting rewrite rules

% Name of the language
% \newcommand{\versecalc}{$\mathcal{VC}$}
% \newcommand{\calcname}{\versecalc\xspace}
%%\newcommand{\mypara}[1]{\smallskip\noindent\textbf{#1}\xspace}  ACMART gives warning about use of \\smallskip for some reason
% \newcommand{\mypara}[1]{\vskip\smallskipamount\noindent\textbf{#1}\xspace}

% \newcommand{\calcname}{Fresh Calculus\xspace}
% \newcommand{\lcalcname}{Lenient Calculus\xspace}
% \newcommand{\calcshort}{FC\xspace}
% \newcommand{\vcshort}{VC\xspace}
\newcommand{\vcname}{Verse Calculus\xspace}

% Right arrow for rewrites
\newcommand{\movesto}{\longrightarrow}
\newcommand{\movesfrom}{\longleftarrow}
\newcommand{\convertible}{\longleftrightarrow^{\ast}}
\newcommand{\movestoX}{\rightsquigarrow}
\newcommand{\movesfromX}{\leftsquigarrow}

% Arrows with explanation, including the common cases of one, two or a repeated rule
\newcommand{\movestoC}[1]{\movesto\!\!\{#1\}}
\newcommand{\movestoR}[1]{\movestoC{\rulename{#1}}}
\newcommand{\movestoRR}[2]{\movestoC{\rulename{#1}, \rulename{#2}}}
\newcommand{\movestoRn}[2]{\movestoC{\rulename{#1} \times #2}}


\newcommand{\rightarrowdbl}{\rightarrow\mathrel{\mkern-14mu}\rightarrow}
\newcommand{\xrightarrowdbl}[2][]{%
  \xrightarrow[#1]{#2}\mathrel{\mkern-14mu}\rightarrow
}
\newcommand{\leftarrowdbl}{\leftarrow\mathrel{\mkern-14mu}\leftarrow}
\newcommand{\xleftarrowdbl}[2][]{%
  \leftarrow\mathrel{\mkern-14mu}\xleftarrow[#1]{#2}%
}

\newcommand{\xleftrightarrowdbl}[2][]{%
  \leftarrow\mathrel{\mkern-14mu}\leftarrow\mathrel{\mkern-14mu}\xrightarrow[#1]{#2}\mathrel{\mkern-14mu}\rightarrow}


\newcommand{\rr}{\ensuremath{\mathcal{R}}\xspace}
\newcommand{\rs}{\ensuremath{\mathcal{S}}\xspace}
\newcommand{\rt}{\ensuremath{\mathcal{T}}\xspace}
\newcommand{\rrhat}{\trel{\rr}}
\newcommand{\rshat}{\trel{\rs}}
\newcommand{\rthat}{\trel{\rt}}
\newcommand{\goone}[1]{\xrightarrow{\epsilon}_{\alignsubscript{#1}}}
\newcommand{\go}[1]{\xrightarrow{}_{\alignsubscript{#1}}}
\newcommand{\gos}[1]{\xrightarrowdbl{}_{\alignsubscript{#1}}}
\newcommand{\lgos}[1]{\xleftarrowdbl{}_{\alignsubscript{#1}}}
\newcommand{\gok}[2]{\xrightarrowdbl{#2}_{\alignsubscript{#1}}}

\newcommand{\trel}[1]{\ensuremath{\widehat{#1}}}
\newcommand{\tstep}[3]{\ensuremath{#2 \go{\trel{#1}} #3}}

\newcommand{\rstepo}[3]{\ensuremath{#2 \goone{#1} #3}}
\newcommand{\rstep}[3]{\ensuremath{#2 \go{#1} #3}}
\newcommand{\rsteps}[3]{\ensuremath{#2 \gos{#1} #3}}
\newcommand{\lsteps}[3]{\ensuremath{#2 \lgos{#1} #3}}
\newcommand{\rstepk}[4]{\ensuremath{#3 \gok{#1}{#2} #4}}
\newcommand{\ungoone}[1]{\xleftarrow{\epsilon}_{\alignsubscript{#1}}}
\newcommand{\ungo}[1]{\xleftarrow{}_{\alignsubscript{#1}}}
\newcommand{\ungos}[1]{\xleftarrowdbl{}_{\alignsubscript{#1}}}
\newcommand{\ungok}[2]{\xleftarrowdbl{#2}_{\alignsubscript{#1}}}
\newcommand{\unrstepo}[3]{\ensuremath{#2 \ungoone{#1} #3}}
\newcommand{\unrstep}[3]{\ensuremath{#2 \ungo{#1} #3}}
\newcommand{\unrsteps}[3]{\ensuremath{#2 \ungos{#1} #3}}
\newcommand{\unrstepk}[4]{\ensuremath{#3 \ungok{#1}{#2} #4}}
\newcommand{\univ}{\ensuremath{\mathcal{E}}}
\newcommand{\join}[3]{{#2} \downarrow_{\alignsubscript{#1}} {#3}}
\newcommand{\dia}{\ensuremath{\diamond}}

\newcommand{\rnk}{\ensuremath{\rho}}
\newcommand{\rord}{\ensuremath{\mathcal{O}}}
\newcommand{\runi}{\ensuremath{\mathcal{U}}}
\newcommand{\runihat}{\ensuremath{\trel{\mathcal{U}}}}
\newcommand{\rstr}{\ensuremath{\mathcal{S}}}
\newcommand{\rapp}{\ensuremath{\mathcal{A}}}
\newcommand{\rapphat}{\ensuremath{\trel{\mathcal{A}}}}
\newcommand{\rcon}{\ensuremath{\mathcal{N}}}
\newcommand{\rconhat}{\ensuremath{\trel{\mathcal{N}}}}
\newcommand{\rconm}{\ensuremath{\rcon'}}
\newcommand{\rss}{\ensuremath{\mathcal{SS}}}

\newcommand{\rgc}{\ensuremath{\mathcal{G}}}
\newcommand{\rgchat}{\ensuremath{\trel{\mathcal{G}}}}
\newcommand{\rfail}{\ensuremath{\mathcal{F}}}
\newcommand{\rspc}{\ensuremath{\mathcal{C}}}
\newcommand{\rspchat}{\ensuremath{\trel{\mathcal{C}}}}
\newcommand{\norecursion}{no recursion\xspace}
\newcommand{\noHYPHENrecursion}{no-recursion\xspace}
\newcommand{\rdx}[1]{\ensuremath{\Delta_{\alignsubscript{#1}}}}
\newcommand{\txr}[2]{\ensuremath{\rstep{#1}{\rdx{#2}}{\rdx{#2}'}}}
\newcommand{\tx}[1]{\ensuremath{\rstep{#1}{\rdx{#1}}{\rdx{#1}'}}}
\newcommand{\txi}[1]{\ensuremath{\rstep{}{\rdx{#1}}{\rdx{#1}'}}}


\newcommand{\wf}[1]{\ensuremath{\vdash #1}}
\newcommand{\wff}[5]{\ensuremath{#1 \vdash #2 \rightsquigarrow (#3 \mid #4 \mid #5)}}
\newcommand{\wffv}[4]{\ensuremath{\vdash_\ensuremath{\Varid{v}} #1 \rightsquigarrow (#2 \mid #3 \mid #4)}}
\newcommand{\wffmany}[4]{\ensuremath{\vdash_{\!*} #1 \rightsquigarrow (#2 \mid #3 \mid #4)}}
% \newcommand{\lvc}[1]{\colorbox{lightgray}{#1}}
% \newcommand{\lvc}[1]{\colorbox{lightgray}{#1}}
\newcommand{\lvc}[1]{#1}
\newcommand{\relDescription}[1]{\ensuremath{\textrm{\textbf{#1}}}}
\newcommand{\judgementHead}[2]{\ensuremath{\relDescription{#1}\hfill\fbox{#2}}}
\newcommand{\judgementHeadNameOnly}[1]{\ensuremath{\relDescription{#1}\hfill}}


%% Definitions for explaining skew confluence
\newcommand{\inccon}[2]{\mathbin{\smash{\mathop{\xleftrightarrowdbl{\hphantom{#1}}}\limits_{\smash{\hbox{\raise0.55ex\hbox{\smash{$\scriptstyle#2$}}}}}^{\vbox to 0pt{\hbox{\lower0.95ex\hbox{\smash{$\scriptstyle#1$}}}\vss}}}}}
\newcommand{\leqomega}{\mathrel{\leq_\omega}}
\newcommand{\omegarr}{\omega_{\rr}}
\newcommand{\leqomegarr}{\mathrel{\leq_{\omegarr}}}
\newcommand{\preceqomegarr}{\mathrel{\preceq_{\omega_{\rr}}}}
\newcommand{\downclose}{\mathord{\Downarrow_{\leqomega}\!}}


%% GS: Add @@{} to match other changes made below to get more horizontal space in Figure 3
%% \newcommand{\rulegroupname}[1]{\multicolumn{4}{l}{\textit{#1}}\\}
\newcommand{\rulegroupname}[1]{\multicolumn{4}{@{}l}{\textit{#1}}\\}
\newcommand{\largerhs}[1]{\multicolumn{2}{@{}l}{#1}}



\newcommand{\ie}{\emph{i.e.}}   % \mbox keeps the last period from
\newcommand{\eg}{\emph{e.g.}}
\newcommand{\Eg}{\emph{E.g.}}
\newcommand{\etc}{\emph{etc.}}
\newcommand{\etal}{\emph{et al.}}
\newcommand{\vs}{\emph{vs.}}

\newcommand{\verifier}{\textsc{TRSVerifier}\xspace}
\newcommand{\mypara}[1]{\vskip\smallskipamount\noindent\textbf{#1}\xspace}
%%%%%%%%%%%%%%%%%%%%%%%%%%%%%%%%%%%%%%%%%%%%%%%%%%%%%%%%%%%%%%%%%%%%%%%%
\newcommand{\vcshort}{VC\xspace}
\newcommand{\isdone}[1]{\mathsf{Verified}(#1)}
\newcommand{\mustKW}{\mathsf{Must}}
\newcommand{\must}[2]{\ensuremath{\mustKW{}\langle #1 \rangle\{#2\}}}
\newcommand{\mustSucceed}[1]{\must{\ensuremath{\textbf{succeeds}}}{#1}}
\newcommand{\mustDecide}[1]{\must{\ensuremath{\textbf{decides}}}{#1}}
\newcommand{\sucdec}[1]{\mathsf{sd}(#1)}
\newcommand{\mustopKW}{\mathsf{Must_{op}}}
\newcommand{\mustop}[2]{\ensuremath{\mustopKW{}\langle #1 \rangle\{#2\}}}

\newcommand{\mustVerifyKW}{\mathsf{Prove}}
\newcommand{\mustVerify}[1]{\ensuremath{\mustVerifyKW(#1)}}
\newcommand{\proves}[2]{\ensuremath{#1 \vdash #2}}
\newcommand{\lamvars}[1]{\mathsf{lams}(#1)}
\newcommand{\UnAssume}[1]{\ensuremath{\lfloor #1 \rfloor}}
\newcommand{\impl}[2]{{#1} \vdash {#2}}
\newcommand{\FEnv}[1]{\ensuremath{\Gamma(#1)}}
\newcommand{\Facts}[1]{\ensuremath{\mathsf{Facts}(#1)}}

\newcommand{\tlamFX}[1]{\ensuremath{\mathsf{tlam}_{#1}}}
% \newcommand{\De}[1]{\llbracket #1 \rrbracket}
% \newcommand{\Fe}[2]{\llbracket #1 \triangleright #2 \rrbracket}
% \newcommand{\Me}[2]{\llbracket #1 \blacktriangleright #2 \rrbracket}
\newcommand{\Ae}[2]{\llbracket #1 \ \textbf{as}\ #2 \rrbracket}


\newcommand{\vb}{\; \ensuremath{\mathop{\mid}} \;}
\newcommand{\textquote}{'}

% Get rid of this lot -------------------
\newcommand{\vtt}[1]{\text{\tt #1}}
\newcommand{\eff}{\mathit{eff}}   % Non-terminal for effects
\newcommand{\clamx}[2]{lam (#1) in #2}
\newcommand{\clam}[2]{\vtt{\clamx{#1}{#2}}}
\newif\ifmatch
\matchtrue
\ifmatch
\newcommand{\pmatch}[2]{$#1 \vtt{ :- } #2$}
\else
\newcommand{\pmatch}[2]{$#1 \leftarrow #2$}
\fi
\newcommand\mydots{\makebox[0.55em][c]{\text{$\cdot\hspace{-0.3em}\cdot\hspace{-0.3em}\cdot$}}}


% ------------------------------------------------------------
%                        Desugaring
% ------------------------------------------------------------

\newcommand{\An}[1]{{\mathcal A}\llbracket #1 \rrbracket}
\newcommand{\Cr}[1]{{\mathcal C}\llbracket #1 \rrbracket}
\newcommand{\Bi}[1]{{\mathcal B}\llbracket #1 \rrbracket}
\newcommand{\Vr}[1]{{\mathcal V}\llbracket #1 \rrbracket}
\newcommand{\Vl}[1]{{\mathcal S}\llbracket #1 \rrbracket}
\newcommand{\Th}[1]{{\mathcal H}\llbracket #1 \rrbracket}

% Control if desugaring has named transformations or not
\newif\ifdesugarops
\desugaropstrue
% \desugaropsfalse

\ifdesugarops
\newcommand{\Ane}[1]{{\mathcal A}\llbracket #1 \rrbracket}
\newcommand{\Fne}[1]{{\mathcal F}\llbracket #1 \rrbracket}
\newcommand{\De}[1]{{\mathcal D}\llbracket #1 \rrbracket}
\newcommand{\Ve}[1]{{\mathcal V}\llbracket #1 \rrbracket}
\newcommand{\Ie}[1]{{\mathcal I}\llbracket #1 \rrbracket}
% \newcommand{\Me}[2]{{\mathcal M}\llbracket #1 \blacktriangleright #2 \rrbracket}
\newcommand{\Me}[2]{{\mathcal M}\llbracket #2 \rrbracket #1}
\newcommand{\DMe}[2]{{\mathcal M}\llbracket #2 \rrbracket #1}
\newcommand{\Fe}[2]{{\mathcal F}\llbracket #2 \rrbracket #1}
\newcommand{\MMe}[3]{{\mathcal M}_{#3}\llbracket #2 \rrbracket #1}
\newcommand{\Ld}[2]{{\mathcal L}\llbracket #1 \rrbracket \, #2}
\newcommand{\Pd}[2]{{\mathcal P}\llbracket #1 \rrbracket \, #2}
\else
\newcommand{\De}[1]{#1}
\newcommand{\Me}[1]{#1}
\newcommand{\Ld}[2]{#1 : #2}
\newcommand{\Pd}[2]{#1 := #2}
\fi

% Names of rules
\newcommand{\rulename}[1]{\textsc{\small #1}}

% Hole in context
\newcommand{\hole}{\square}
\newcommand{\freevars}[1]{\mathsf{fvs}(#1)}
\newcommand{\skol}[1]{\mathsf{skol}(#1)}
\newcommand{\freevarsoutsidelambdas}[1]{\mathsf{fvsol}(#1)}
\newcommand{\boundvars}[1]{\mathsf{bvs}(#1)}
\newcommand{\valvars}[1]{\mathsf{vvs}(#1)}
\newcommand{\disjoint}{\;\#\;}
\newcommand{\hnf}{\mathit{hnf}}  % Head normal form

\newcommand{\hexp}{h-expression}
\newcommand{\Hexp}{H-expression}

\newcommand{\coloredtext}[3]{\strut\ifmmode\text{\color{#1} [{#2}: #3]}\else{\color{#1} [{#2}: #3]}\fi}
% \text in text mode implies a mbox and prevents wrapping, so only use \text in math mode
% The strut is needed somehow, else \iffmode does not work in an \begin{array} cell?
\definecolor{dkcyan}{rgb}{0.1, 0.3, 0.3}
\newcommand{\asktim}{\text{\color{red} $\ast$}}
\newcommand{\lennart}{\coloredtext{dkcyan}{LA}}
\newcommand{\simon}{\coloredtext{red}{SPJ}}
\newcommand{\ranjit}{\coloredtext{green}{RJ}}
\definecolor{dkblue}{rgb}{0.0, 0.0, 0.6}
\newcommand{\joachim}{\coloredtext{dkblue}{JB}}

\newenvironment{figurebox}{\begin{figure}\begin{framed}}{\end{framed}\end{figure}}
\newenvironment{figurebox*}{\begin{figure*}\begin{framed}}{\end{framed}\end{figure*}}

% Example environment
\newcounter{example}
\newenvironment{example}[1][]{\refstepcounter{example}\par\medskip
   \noindent\textbf{Example~\theexample. #1} \rmfamily}{\medskip}


% ---------------------------------------------------
%             Denotational semantics
% ---------------------------------------------------

\newcommand{\M}[1]{\ensuremath{#1}\xspace}
\newcommand{\SP}{\hspace{.5mm}}
\newcommand{\SPP}{\hspace{1mm}}

%% Functions
\newcommand{\parfun}[2]{\M{#1 \rightharpoonup #2}} % Partial function
\newcommand{\totfun}[2]{\M{#1 \rightarrow #2}}     % Total function
\newcommand{\semfcn}{\parfun{\Value}{\Value}}      % Partial function W->W
\newcommand{\dom}{\mathbf{dom}}

%% Domains
\newcommand{\Term}{Term}
\newcommand{\Labels}{Labels}
\newcommand{\Forml}{Form}
\newcommand{\Env}{\M{\mathit{Env}}}
\newcommand{\ENV}{\M{\mathit{ENV}}}
\newcommand{\Ident}{\M{\mathit{Iden}}}
\newcommand{\Fun}{\M{\mathit{Fun}}}
\newcommand{\Z}{\M{\mathbb Z}}
\newcommand{\Val}{\M{W}}
\newcommand{\Value}{\M{W}}

%% Sets and sequences
\newcommand{\powerset}[1]{\M{\mathit{Set}(#1)}}    % Sets
\newcommand{\seq}[1]{\M{ [\SP #1 \SP] }}           % Sequences

%% Semantic function calls, e.g E[e]rho
\newcommand{\semfun}[1]{\mathcal{#1}}
\newcommand{\semap}[2]{\semfun{#1}\llbracket #2 \rrbracket}
\newcommand{\semapp}[3]{\semfun{#1}\llbracket #2 \rrbracket\, #3}
\newcommand{\semappp}[4]{\semfun{#1}\llbracket #2 \rrbracket\, #3 \, #4}
\newcommand{\semapppp}[5]{\semfun{#1}\llbracket #2 \rrbracket\, #3 \, #4 \, #5}

% Functions over ENV
\newcommand{\EV}{\Delta}   % A variable that ranges over \ENV
\newcommand{\univE}{\mathbf{univ}}
\newcommand{\equalsE}{\approx}
\newcommand{\cleanE}{\mathbf{clean}}
\newcommand{\intersectE}{\cap}
\newcommand{\unionE}{\cup}
\newcommand{\hideE}{\mathbf{hide}}

% Functions over sequences
\newcommand{\append}{+\!\!+\,}

\newcommand{\tuple}[1]{\langle #1 \rangle}
\newcommand{\equa}[1]{\begin{equation*}#1\end{equation*}}
\newcommand{\mcompr}[2]{[ #1 \vb #2 ]}
\newcommand{\mbind}[2]{ #1 \leftarrow #2 }
\newcommand{\malt}{\diamond}
\newcommand{\semclose}[2]{\textbf{close}_{#1}\,#2}
\newcommand{\semopen}[2]{\textbf{open}_{#1}\,#2}
\newcommand{\semapply}[2]{\apply(#1,#2)}
\newcommand{\apply}{\mathbf{apply}}

\newcommand{\setform}[1]{\left\{ #1 \right\}}
\newcommand{\lcompr}[2]{[ #1 \vb #2 ]}
\newcommand{\scompr}[2]{\{ #1 \vb #2 \}}
\newcommand{\sbind}[2]{ #1 \in #2 }
\newcommand{\inb}{\,\dot{\in}\,}

\newcommand{\envapp}[2]{\semfun{R}\llbracket #1 \rrbracket #2}
\newcommand{\folapp}[2]{\semfun{L}\llbracket #1 \rrbracket #2}
\newcommand{\sfolapp}[2]{\semfun{S}\llbracket #1 \rrbracket #2}
\newcommand{\termapp}[2]{\semfun{T}\llbracket #1 \rrbracket #2}
\newcommand{\liftrel}[1]{\hat{#1}}
\newcommand{\cons}{{:}}           % Cons for value sequences, V*

\newcommand{\sconcat}{\bowtie}
\newcommand{\shead}{\mathit{head}}
\newcommand{\stuple}{\mathit{tuple}}
\newcommand{\scons}{:}

\newcommand{\fcn}{\mathbf{fcn}}
\newcommand{\get}[1]{\mathbf{sing}(#1)}

\newcommand{\dtuple}{\tuple{\Value}}
\newcommand{\fmapsto}{\!\mapsto\!}
\newcommand{\mkfcn}[1]{\fcn #1}


% extend envt
\newcommand{\extendrho}[2]{#1[#2]}

% The collection type
\newcommand{\mtype}{\Value^{*}}
\newcommand{\mtyped}{\Value^{?}}
\newcommand{\mempty}{\{\}}
\newcommand{\munit}{\textbf{unit}}
\newcommand{\msing}[1]{\{#1\}}
\newcommand{\munits}[1]{\munit(#1)}
\newcommand{\mintersect}{\Cap}
\newcommand{\munion}{\Cup}
\newcommand{\msemi}{\fatsemi}
\newcommand{\mone}{\mathit{one}}
\newcommand{\mall}{\mathit{all}}
\newcommand{\msucceeds}{\mathit{succeeds}}
\newcommand{\mdecides}{\mathit{decides}}
\newcommand{\mtag}{\mathit{tag}}
\newcommand{\mtage}[1]{([],#1)}
\newcommand{\mlift}[1]{#1_{\bot}}
\newcommand{\mplus}{\oplus}
\newcommand{\mcross}{\scalebox{1.5}{$\Cup$}}
\newcommand{\wrong}{\mathbf{wrong}}
\newcommand{\ran}{\mathbf{ran}}
